\begin{abstract}
Vision-language models usually reason about change in text. Would they do better if they could \emph{see} the changed scene? We build a benchmark in which the state a question asks about (the view from another position, the scene after two objects trade places, or the outcome of releasing a stack of blocks) can be rendered exactly, and give it to the model as a second image. Across three open VLMs from two families, a perfect imagined state raises accuracy on every task, by up to 30 points, while describing the change in text first gives a significant gain in only one case, a question that can be restated in terms of the visible scene. On viewpoint questions, a misleading imagined state drops accuracy far below the image-only baseline, and models then give the answer it implies up to 81\% of the time, so the visual state plays a causal role. Imagination need not be precise: an oracle with 1\,m of position noise still helps on viewpoint questions. Yet models gain little from their own imagined states, partly because they misperceive the relations that matter and partly because they do not rely on what they produced. Finally, the seen image competes with the imagined one: on viewpoint questions whose answer changes after moving, models answer 90--98\% correctly from the imagined view alone but only 29--61\% when the original image is also shown.
\end{abstract}
